% ============================================================================
% IV. WHEN DOES JOINT CONTROL PAY?  (theory)
% Formal content: sections/theory_formal.tex (from docs/research_v2/
% theory_latex.tex, corrected version merged at 1d8900c). Source of truth: docs/research_v2/
% THEORY.md. This file holds the author-prose outline around it and the
% prediction table that the v2 test is built on.
% Framing (PI decision, 2026-10-15): the paper is a boundary-condition
% result. Equivalence under A1--A3; C1 is second order; joint control matters
% mainly when a dispatched retrieval is COMMITTED; under SPECULATIVE prefetch
% a belief-gated decomposed prefetcher (B11, SPEC-theta) is the strong rival
% (Prop. 4''). v2 runs a pre-registered boundary test (P1, P2, P3).
% ============================================================================
\section{When Does Joint Control Pay?}\label{sec:theory}

% ---- P-IV.1 -----------------------------------------------------------------
% PURPOSE: state assumptions A1--A3 in words and say that v1 satisfied all
%   three BY DESIGN, so the v1 tie was predictable (R1.1, R2 "keep the
%   coupling null central").
% CLAIMS + EVIDENCE:
%   (A1) quality independent of timing: v1 retrieved on a frozen,
%        precomputed query, so partial-query risk was structurally zero
%        [simulator/e2_results/RUN_MANIFEST.md; _objective.py
%        partial_query_risk]; E3 answers were generated once per (scenario,
%        seed, path) and reused for every timing [simulator/e3_full_sweep.py].
%   (A2) constant overlap per policy: v1 used one overlap per policy
%        (0, 0.050, 0.188 s) [V1T constants.overlap_s].
%   (A3) one belief per turn computed from the full question
%        [simulator/e3_scenario_signals.json].
\todo{P-IV.1: the three assumptions and why v1 met each by construction.}

% ---- P-IV.2 -----------------------------------------------------------------
% PURPOSE: Prop. 1 and Cor. 1 in words; the exact price map; B9+ (R1.1:
%   "a decomposed controller that prices expected early-dispatch latency
%   equals Ours") -- concede the point and make it a result.
% CLAIMS: Ours at lambda_l = B9 at lambda_l' (exact, including the waste
%   term); B9+ = Ours at the same price (Prop. 2, last line).
% NUMBERS: none here (they belong to Sec. VI).
\todo{P-IV.2: Proposition~\ref{prop:eq} and its corollary in words; the
decomposed controller that prices expected exposure (\Bplus) is the joint
controller under (A1)--(A3).}

% ---- P-IV.3 -----------------------------------------------------------------
% PURPOSE: Prop. 2 -- the divergence band and the gap bound
%   lambda_l * o* * P(b in band) (R1.1 "derive the divergence band").
% CLAIMS: the band is B_lambda = (theta_O, theta_B]; outside it the two
%   decisions coincide; the gap is bounded by lambda_l o* P(band) under model
%   correctness; without model correctness it can be negative (Remark).
\todo{P-IV.3: Proposition~\ref{prop:gap}, the band and the bound, and why the
bound is a statement about the controller's model, not about outcomes.}

\begin{equation}
\mathcal B_{\lambda}=\left(\theta_O,\theta_B\right],\quad
\theta_O=\frac{c_O+\gm}{s},\quad \theta_B=\frac{\lamL L+\lamC(1+w_{\mathrm{se}})+\gm}{s},
\label{eq:band}
\end{equation}
with $c_O=\min_t[\lamL e_t+\lamC(1+w_t)]$ and $s=\gp(1-\bar\rho)+\gm$; for
the v1 objective $\mathcal B_\lambda=(\lamL\cdot0.484+1.08\,\lamC,\ \lamL\cdot0.672+\lamC]$ at
the maximal-overlap timing.

% ---- P-IV.4 -----------------------------------------------------------------
% PURPOSE: the boundaries -- which channel breaks the equivalence, and how
%   much (R2: "state where joint control could help").
% CLAIMS + EVIDENCE [TH Secs. 5--7]:
%   (a) C1, turn-dependent overlap (drop A2): gain is a conditional Jensen gap
%       and SECOND ORDER, about f(theta*) lambda_l^2 Var(e*) / (2(g^+ + g^-))
%       (Prop. 3). In simulation the best cell needs about 20,821 test items
%       for 80% power against B9+ [POW; PRE2 Sec. 6] -- C1 cannot carry a
%       confirmatory test at the planned sizes.
%   (b) C2, information versus overlap (drop A3): waiting buys a sharper
%       belief (Howard's value of information [howard1966information]) and
%       loses overlap; under COMMIT the optimal stopping region is not a
%       product set, so every decomposed policy is strictly suboptimal when
%       both regions have mass (Prop. 4).
%   (c) Under SPECULATIVE prefetch (Prop. 4'): the joint controller's exact
%       gain over always-prefetch B10 is the expected positive part of the
%       value of waiting over the prefetch option. It is bounded by the
%       compute spent on turns where joint does not prefetch plus a C3 term,
%       so B10 is optimal when lambda_c = 0 and Delta rho = 0. The margin
%       lambda_c(1+w)q is NOT a bound (counterexample in the proposition);
%       the valid two-epoch bound is m_prophet.
%   (e) Prop. 4'' -- WHY SPEC-THETA MATTERS (outline note for the authors):
%       the joint controller beats B10 mainly by declining to prefetch on
%       confident-skip prefixes. That decision uses only the belief at the
%       trigger and fixed prices, so a decomposed controller can make it.
%       B11 makes it, and its gap to the joint optimum is bounded by
%       lambda_l E|X_hat - X_bar|(u* + u_theta), a C1-sized term (second
%       order in expectation). Under speculation the honest decomposed
%       rival is therefore B11; comparing joint with B10 alone would credit
%       joint control with a decomposable compute saving, and P2 would then
%       be predicted false. In simulation, joint minus B11 lies in
%       [-0.0005, +0.0009] across 24 cells, while joint minus B10 is 0.012
%       to 0.023 above the median cut [TH Sec. 6.2; not results].
%   (f) What can still keep joint ahead under speculation: delaying a paid
%       prefetch across several partials, and C3 query improvement; the
%       two-epoch simulation does not quantify either [TH Sec. 6.2].
%   (d) C3, partial-query risk alone (drop A1): optimal timing becomes
%       belief-dependent but stays b-measurable, so it does not separate the
%       classes (Prop. 5).
% NUMBERS: 20,821 [POW cells_C1 / PRE2 Sec. 6]; [-0.0005, +0.0009] and
%   0.012-0.023 [TH Sec. 6.2, simulation, not results].
\todo{P-IV.4: the four channels and their sizes; joint control is predicted
to matter when retrieval is committed; under speculation a belief-gated
decomposed prefetcher (\SPECth) is predicted to match it, which is why it is
a required rival (Propositions~\ref{prop:spec} and~\ref{prop:spectheta},
Table~\ref{tab:predictions}).}

% ============================================================================
% THEORY SLOT -- verbatim copy of docs/research_v2/theory_latex.tex as merged
% into research/v2-base at 1d8900c (corrected boundary test: Prop. 4' exact
% gain and bound for SPEC (B10), Prop. 4'' for the belief-gated prefetcher
% SPEC-theta (B11)). Source of truth: docs/research_v2/THEORY.md; numbers in
% docs/research_v2/v1_theory_check.json and v2_power.json (block
% boundary_test). If THEORY.md changes, re-copy this file; do not edit here.
%
% AUTHOR ACTIONS:
%  * Statements and equations may stay. The connecting sentences and proof
%    wording must be reviewed and re-worded by the authors before submission
%    (conference AI-content rule).
%  * Page budget: the full proofs will not fit in a 6-page INDICON paper.
%    Keep the statements and move the proofs to a supplementary file or an
%    extended version: \PEND{decide proof placement}.
%  * Any "Verification on the v1 sweep" subsection duplicates the v1 results
%    section (label sec:v1); keep it in one place only.
% ============================================================================
% ============================================================================
% theory_latex.tex -- propositions and proofs for the revised paper.
% Source of truth: docs/research_v2/THEORY.md (numbers: v1_theory_check.json,
% v2_power.json). Paste into a section, or \input{} it.
%
% Preamble requirements (IEEEtran-compatible). main.tex currently loads
% amsmath/amssymb/amsfonts but NOT amsthm; add to the preamble:
%   \usepackage{amsthm}
% The theorem environments below are defined only if absent. If the class's
% own proof environment clashes, replace proof with IEEEproof.
% Not compile-tested in this checkout (no TeX installation); compile once in
% the paper build before submission.
% ============================================================================
\ifdefined\proposition\else\newtheorem{proposition}{Proposition}\fi
\ifdefined\lemma\else\newtheorem{lemma}{Lemma}\fi
\ifdefined\corollary\else\newtheorem{corollary}{Corollary}\fi
\ifdefined\remark\else\newtheorem{remark}{Remark}\fi

\subsection{Model}
A turn has a need indicator $N\in\{0,1\}$ and a calibrated belief
$b=\Pr(N=1\mid\mathcal I)$. Retrieval changes correctness by
$g^{+}=\mathbb E[Q^{r}-Q^{s}\mid N=1]$ and $-g^{-}=\mathbb E[Q^{r}-Q^{s}\mid N=0]$.
A timing action $t\in\mathcal T$ hides overlap $o_t$ of a retrieval of latency
$L$, so it exposes $e_t=L-o_t$. It wastes $w_t$ extra retrievals and loses a
fraction $\rho_t$ of $g^{+}$ to a partial query. With prices
$\lambda_l,\lambda_c\ge 0$, $U(\text{skip})=0$ and
\begin{equation}
U_t(b)=b\,g^{+}(1-\rho_t)-(1-b)\,g^{-}-\kappa_t,\quad
\kappa_t=\lambda_l e_t+\lambda_c(1+w_t).
\label{eq:U}
\end{equation}
The v1 objective is $g^{+}=1$, $g^{-}=\rho_t=0$. The joint controller (Ours)
maximizes \eqref{eq:U} over skip and all $(\text{retrieve},t)$, with ties to
skip. The decomposed controller (B9) retrieves iff
$s b-g^{-}-\lambda_lL-\lambda_c(1+w_{\mathrm{se}})>0$, pricing the full
latency. Its timing, chosen independently, is
$\arg\max_t[-bg^{+}\rho_t+\lambda_lo_t-\lambda_cw_t]$. B9+ uses the same
timing $\hat\tau$ but prices its expected exposure: it retrieves iff
$U_{\hat\tau}(b)>0$.

The policy class $\Pi_J$ allows the retrieve decision $r$ to depend on
$(b,Z)$, where $Z$ is the turn's timing state. The decomposed class
$\Pi_D\subset\Pi_J$ requires $r$ to be $\sigma(b)$-measurable.

The v1 regime assumes:
\textbf{(A1)} $\rho_t\equiv\bar\rho$;
\textbf{(A2)} $e_t$ is constant across turns;
\textbf{(A3)} one belief per turn, independent of dispatch time.
Under (A1), every slope is $s=g^{+}(1-\bar\rho)+g^{-}$.

\begin{lemma}\label{lem:J}
For integrable $X$,
$\mathbb E X^{+}-(\mathbb E X)^{+}=\min(\mathbb EX^{+},\mathbb EX^{-})
=\tfrac12(\mathbb E|X|-|\mathbb EX|)\le\tfrac12\mathbb E|X-\mathbb EX|$.
The gap is zero iff $X\ge0$ a.s.\ or $X\le0$ a.s.
\end{lemma}
\begin{proof}
If $\mathbb EX\ge0$, then $\mathbb EX^{+}-\mathbb EX=\mathbb EX^{-}\le\mathbb EX^{+}$.
The case $\mathbb EX<0$ is symmetric. Use $x^{+}=(x+|x|)/2$ and
$|X|\le|X-\mathbb EX|+|\mathbb EX|$.
\end{proof}

\begin{lemma}\label{lem:I}
For square-integrable $C$,
$\int_{\mathbb R}\big(\mathbb E(u-C)^{+}-(u-\mathbb EC)^{+}\big)\,du=\tfrac12\operatorname{Var}C$.
\end{lemma}
\begin{proof}
Let $m=\mathbb EC$. For large $M$,
$\int_{-M}^{M}[(u-c)^{+}-(u-m)^{+}]du=(m-c)(M-m)+\tfrac12(m-c)^2$.
Take expectations and let $M\to\infty$.
\end{proof}

\subsection{Equivalence and the fixed-price gap}
Let $t^{*}=\arg\min_t\kappa_t$ and $c_O=\kappa_{t^{*}}$, and define
$\lambda_l'=[c_O-\lambda_c(1+w_{\mathrm{se}})]/L$.

\begin{proposition}[Equivalence]\label{prop:eq}
Assume (A1)--(A3) and $w_{\mathrm{se}}=\min_tw_t$. For every $b\in[0,1]$, Ours
at $(\lambda_l,\lambda_c)$ retrieves iff B9 at $(\lambda_l',\lambda_c)$
retrieves. The map $\lambda_l\mapsto\lambda_l'$ is a continuous, strictly
increasing bijection of $[0,\infty)$. When the maximal-overlap timing is optimal,
\[
\lambda_l'=\lambda_l\frac{L-o^{*}}{L}+\lambda_c\frac{w^{*}-w_{\mathrm{se}}}{L},
\]
which reduces to $\lambda_l(L-o^{*})/L$ iff $\lambda_c(w^{*}-w_{\mathrm{se}})=0$.
Conditional on retrieving, the timings coincide unless $\lambda_l$ and
$\lambda_l'$ lie in different linear pieces of $\lambda\mapsto\min_t\kappa_t$.
\end{proposition}
\begin{proof}
Under (A1), $\max_tU_t(b)>0$ iff $sb-g^{-}>c_O$. Ties go to skip, so Ours
retrieves iff $b>\theta_O=(c_O+g^{-})/s$. B9 retrieves iff
$b>\theta_B(\lambda')=(\lambda'L+\lambda_c(1+w_{\mathrm{se}})+g^{-})/s$, and
$\theta_B(\lambda_l')=\theta_O$ by construction. Both comparisons are strict.

The function $\min_t\kappa_t(\lambda)$ is a minimum of affine maps with slopes
$e_t>0$, equal to $\lambda_c(1+w_{\mathrm{se}})$ at $\lambda=0$; hence the
bijection.

Under (A1) the belief term of B9's timing objective is constant in $t$, so its
timing is $\arg\min_t\kappa_t(\lambda')=t^{*}(\lambda')$, whereas Ours uses
$t^{*}(\lambda)$.
\end{proof}

\begin{corollary}[Frontiers]
If each turn's outcome depends only on (item, seed, retrieve, timing), then
sweeping Ours and B9 over $\lambda_l\in[0,\infty)$ yields identical sets of
operating points, except on the timing-mismatch interval of
Proposition~\ref{prop:eq}.
\end{corollary}

\begin{proposition}[Gap at a fixed price]\label{prop:gap}
Assume (A1)--(A3) and model correctness: $\mathbb E[N\mid b]=b$, and realized
quality, latency and compute in expectation equal to those in \eqref{eq:U}.
Let $c_B=\lambda_lL+\lambda_c(1+w_{\mathrm{se}})$ and
$\mathcal B=(\theta_O,\theta_B]$. Then
\[
\mathbb E[U_{\mathrm{Ours}}-U_{\mathrm{B9}}]=\mathbb E[(sb-g^{-}-c_O)\mathbf 1\{b\in\mathcal B\}]
\in\big[0,\ (c_B-c_O)\Pr(b\in\mathcal B)\big],
\]
and $(c_B-c_O)\Pr(b\in\mathcal B)\le\lambda_lo^{*}\Pr(b\in\mathcal B)$.
For B9+ at the same price the difference is $0$ pointwise.
\end{proposition}
\begin{proof}
At equal prices B9's timing is $t^{*}$, so whenever B9 retrieves it pays
$c_O$. The two controllers differ only on $\mathcal B$, where Ours earns
$sb-g^{-}-c_O\in(0,c_B-c_O]$ and B9 earns $0$. B9+ applies Ours' rule
$U_{t^{*}}(b)>0$.
\end{proof}

\begin{remark}
Without model correctness, the realized gap is
$\mathbb E[(m(b)-\tilde c)\mathbf 1\{b\in\mathcal B\}]$, where
$m(b)=\mathbb E[Q^{r}-Q^{s}\mid b]$ and $\tilde c$ is the realized cost. It may
be negative, and it is not $O(\lambda_lo^{*})$. On the v1 sweep the
controller model predicts a positive gap at all 11 band prices. The realized
gap, computed with measured TTFA, is negative at 10 of 11. The prediction made
with the corrected objective and the measured cost agrees in sign at 10 of 11
($r=0.952$).
\end{remark}

\subsection{Channels that break the equivalence}
\textbf{C1 (turn-dependent overlap).} Drop (A2), and let
$e_t(Z)=\mathbb E[\pi+\max(0,L-R_t)\mid Z]$. Write $A(b)=sb-g^{-}$,
$C(Z)=\min_t[\lambda_le_t(Z)+\lambda_c(1+w_t)]$ and $\Delta=A(b)-C(Z)$.

\begin{proposition}[C1]\label{prop:c1}
Under model correctness:
\begin{itemize}
\item $\sup_{\Pi_J}\mathbb EU=\mathbb E\Delta^{+}$ and
$\sup_{\Pi_D}\mathbb EU=\mathbb E(\mathbb E[\Delta\mid b])^{+}$; if $b\perp Z$,
the latter is attained by B9+ pricing $\mathbb EC$.
\item The gain is
$G=\mathbb E_b\big[\min(\mathbb E[\Delta^{+}\mid b],\mathbb E[\Delta^{-}\mid b])\big]\ge0$.
It is strictly positive iff, with positive probability over $b$, $\Delta$
takes both signs given $b$.
\item If $b\perp Z$:
$G\le\tfrac12\mathbb E|C-\mathbb EC|\,\Pr(A(b)\in(\inf C,\sup C))$.
\item If moreover $b$ has a density $f$ that is continuous at
$\theta^{*}=(\mathbb EC+g^{-})/s$, then
$G=\frac{f(\theta^{*})}{2s}\operatorname{Var}C\,(1+o(1))$ as the spread of
$C$ vanishes.
\end{itemize}
\end{proposition}
\begin{proof}
Optimize pointwise in $(b,Z)$ for $\Pi_J$. For $\Pi_D$, condition on $b$:
$\mathbb E[r(b)\Delta]=\mathbb E[r(b)\mathbb E[\Delta\mid b]]$. Subtract the two
values and apply Lemma~\ref{lem:J} conditionally on $b$. For the
approximation, substitute $u=A(b)$ in
$\int f(b)[\mathbb E(A(b)-C)^{+}-(A(b)-\mathbb EC)^{+}]db$ and apply
Lemma~\ref{lem:I}.
\end{proof}

The gain is second order: of order $\lambda_l^{2}\operatorname{Var}e^{*}(Z)$.
If the plug-in cost satisfies $|\hat C-C|\le\varepsilon$, the joint rule loses
at most $\varepsilon\Pr(|\Delta|\le\varepsilon)$, so a realized advantage
requires $G>\varepsilon\Pr(|\Delta|\le\varepsilon)$.

\textbf{C2 (information vs.\ overlap).} Drop (A3). The epochs are
$k=1,\dots,K$, and the belief $b_k=\Pr(N=1\mid\mathcal F_k)$ is a martingale.
A retrieval dispatched at $k$ is committed. Its value is
$D_k=s_kb_k-g^{-}-\lambda_lX_k-\lambda_c(1+w_k)$, and the problem's value is
$V_K=D_K^{+}$, $V_k=\max(D_k,\mathbb E[V_{k+1}\mid\mathcal F_k])$. Let
$\delta_k=\mathbb E[D_{k+1}-D_k\mid\mathcal F_k]$ and
$I_k=\min(\mathbb E[D_{k+1}^{+}\mid\mathcal F_k],\mathbb E[D_{k+1}^{-}\mid\mathcal F_k])$.

\begin{proposition}[C2]\label{prop:c2}
The one-step value of waiting is
$W_k=\mathbb E[D_{k+1}^{+}\mid\mathcal F_k]-D_k^{+}=I_k+(D_k+\delta_k)^{+}-D_k^{+}$.
\begin{enumerate}
\item If $W_k>0$, waiting strictly beats both dispatching and skipping now.
\item If $\delta_j\le0$ for all $j\ge k$, then
$V_k-D_k\le\mathbb E[\max_{j\ge k}D_j^{-}\mid\mathcal F_k]$. Hence dispatching
now is optimal when $D_k>0$ and $\Pr(\min_{j\ge k}D_j<0\mid\mathcal F_k)=0$, and
strictly optimal if in addition $\delta_k<0$.
\item Suppose some event $E$ in the belief-free timing filtration
$\mathcal G_k$ contains, with positive probability, both turns in the strict-wait
set of (1) and turns in the strict-dispatch set of (2). Then every decomposed
policy is strictly suboptimal. Here a decomposed policy dispatches at a
$\mathcal G$-stopping time and decides with $b$ at that time.
\end{enumerate}
\end{proposition}
\begin{proof}
(1) Waiting one step and then acting myopically is feasible. Apply
Lemma~\ref{lem:J} to $D_{k+1}$ given $\mathcal F_k$, using the martingale
property for $\delta_k$.
(2) $(D_j)_{j\ge k}$ is a supermartingale. A stopped payoff is at most
$D_\sigma+\max_jD_j^{-}$; apply optional stopping.
(3) On $E$ a decomposed policy takes the same epoch-$k$ action (commit or
defer) for every belief. It therefore loses a strictly positive amount on one
of the two subsets.
\end{proof}

\textbf{Speculative architecture.} A retrieval prefetched at epoch $k$ costs
$c_k=\lambda_c(1+w_k)$, which is sunk. At $K$ it is used iff
$Y_k=b_Kg^{+}(1-\rho_k)-(1-b_K)g^{-}-\lambda_l\hat X_k>0$, where $\hat X_k$ is
its expected exposure given the realized hiding window. Without a prefetch,
the controller dispatches at $K$ iff
$U_K=A_K(b_K)-\lambda_l\hat X_K-c_K>0$.

The joint controller is the Snell envelope with stopping payoff
$O_k=\mathbb E[Y_k^{+}\mid\mathcal F_k]-c_k$. SPEC always prefetches at the
first query-ready epoch $k_0$. SPEC-$\theta$ prefetches at $k_0$ iff the
belief-only option value, computed with dev-mean exposures $\bar X$, exceeds
the belief-only value of not prefetching.

Write $c=c_{k_0}$, $Y=Y_{k_0}$, $q=\Pr(Y\le0)$ and
$\Delta\rho=\rho_{k_0}-\min_{k>k_0}\rho_k$. Assume exposures are pathwise
non-decreasing in the dispatch epoch.

\begin{proposition}[Speculative dispatch]\label{prop:spec}
Let $G=V^{*}_S-V_{\mathrm{SPEC}}$, and let $D_J$ be the event that the joint
controller does not prefetch at $k_0$.
\begin{enumerate}
\item $G=\mathbb E[(\mathbb E[V_{k_0+1}\mid\mathcal F_{k_0}]-O_{k_0})^{+}]$.
\item $0\le G\le c\Pr(D_J)+g^{+}\Delta\rho\,\mathbb E[b_K\mathbf 1_{D_J}]$.
In particular $G\le\lambda_c(1+w_{k_0})$ when $\Delta\rho=0$, and SPEC is
optimal when $\lambda_c=\Delta\rho=0$.
\item For two epochs and $\Delta\rho=0$,
\[
G\le cq+\mathbb E[(c-Y)\mathbf 1\{0<Y<c\}]
+\mathbb E[(c-c_K-\lambda_l(\hat X_K-\hat X_{k_0}))^{+}\mathbf 1\{Y>0\}].
\]
The quantity $cq$ alone is not a bound. For example, if $Y\in(0,c)$ a.s.,
$b_K$ is known at $k_0$ and $\lambda_l(\hat X_K-\hat X_{k_0})\ge Y$, then
$q=0$ but $G=c-\mathbb EY>0$.
\end{enumerate}
\end{proposition}
\begin{proof}
(1) SPEC stops at $k_0$, whereas the optimum takes the larger of stopping and
continuing.

(2) Every later option is pathwise at most
$Y^{+}+g^{+}b_K(\rho_{k_0}-\rho_k)$, because exposures are non-decreasing and
$(x+d)^{+}\le x^{+}+d$ for $d\ge0$. Hence
$\mathbb E[V_{k_0+1}\mid\mathcal F_{k_0}]\le O_{k_0}+c+g^{+}\Delta\rho\,\mathbb E[b_K\mid\mathcal F_{k_0}]$.

(3) Compare SPEC with the ex-post best alternative, which is not prefetching.
\end{proof}

\begin{proposition}[Belief-gated prefetch]\label{prop:spectheta}
Assume two epochs and $\Delta\rho=0$. Then
\[
0\le V^{*}_S-V_{\mathrm{SPEC}\text{-}\theta}
\le\lambda_l\,\mathbb E\big[|\hat X_{k_0}-\bar X_{k_0}|(u^{*}+u^{\theta})\big],
\]
where $u$ indicates ``prefetch at $k_0$ and use''. SPEC-$\theta$ is decomposed,
so the joint advantage under speculation is a C1-sized term.
\end{proposition}
\begin{proof}
With the scalar exposure $\bar X_{k_0}$, only the belief matters and the
optimum is SPEC-$\theta$'s gate. Replacing $\hat X_{k_0}$ by $\bar X_{k_0}$
changes any policy's payoff by at most $\lambda_l|\hat X_{k_0}-\bar X_{k_0}|$
on prefetch-and-use paths. SPEC-$\theta$'s realized-window use rule only
improves on its scalar version.
\end{proof}

\textbf{C3 (partial-query risk).} Keep (A2) and (A3), and let $\rho_t$ vary.

\begin{proposition}[C3]\label{prop:c3}
The optimal timing given retrieval, $t^{*}(b)=\arg\max_t[-g^{+}\rho_tb-\kappa_t]$,
is non-decreasing in $b$ when later timings have lower $\rho_t$ and higher
$\kappa_t$. The retrieve set $\{b:\max_tU_t(b)>0\}$ is an upper interval
$(\theta,1]$. Hence the joint optimum is $\sigma(b)$-measurable: C3 alone does
not separate $\Pi_J$ from $\Pi_D$. It defeats only decompositions whose timing
ignores $b$, and scalar-priced rules such as B9+.
\end{proposition}
\begin{proof}
$\partial_bU_t=s_t$ increases in $t$, so $U$ has increasing differences in
$(b,t)$, and Topkis' theorem gives monotonicity of $t^{*}$. Each $U_t$ is
increasing in $b$, so the retrieve set is an upper interval. The pair
$(\mathbf 1\{b>\theta\},t^{*}(b))$ is $b$-measurable and attains the optimum.
\end{proof}

\subsection{Corrected objective}
With $\rho\equiv0$, retrieval at cost $\kappa$ pays iff
$b>b^{*}=(\kappa+g^{-})/(g^{+}+g^{-})$. On v1,
$\hat g^{+}=0.720\ [0.504,0.888]$ and $\hat g^{-}=0.592\ [0.408,0.768]$, so
$b^{*}(0)=0.451$, against $0.05$ under the v1 objective. The correction leaves
Proposition~\ref{prop:eq} intact, since the rule is still a threshold in $b$.
It places the zero-price boundary mid-range.

\subsection{Verification on the v1 sweep}
\begin{itemize}
\item Ours at grid price $\lambda_i$ and B9 at $\lambda_{i-1}$ agree on
1{,}250/1{,}250 retrieve decisions and on 941/941 timings.
\item On 200 off-grid prices, the exact map agrees on 10{,}000/10{,}000 decisions.
The approximate map $\lambda(L-o^{*})/L$ agrees on 9{,}938/10{,}000.
\item The band is non-empty at 11 of 26 prices. It contains 41 of the 1{,}300
decisions (3.15\%), and the observed disagreements equal the in-band set at
every price.
\item B9+ reproduces Ours on all 11{,}300 decisions checked.
\item The grid satisfies
$(\lambda_i-\lambda_{\deg})/(\lambda_{i-1}-\lambda_{\deg})=L/(L-o^{*})=1.3884$,
with $\lambda_{\deg}=0.02128$. This holds by construction.
\end{itemize}


% Table -- pre-registered predictions (complete; predictions, not results).
% Sources: docs/research_v2/THEORY.md Sec. 7; docs/research_v2/PREREG_V2.md
% Sec. 3. The two-architecture rows reflect the PI decision of 2026-10-15
% (test the boundary in both directions); PREREG_V2.md must be updated to
% match before it is frozen.
\begin{table}[t]
\caption{Pre-Registered Directional Predictions for v2. $\bar D$ is the mean
paired per-item utility difference (Sec.~\ref{sec:setup}).}
\label{tab:predictions}
\centering
\footnotesize
\setlength{\tabcolsep}{3pt}
\begin{tabular}{@{}p{2.45cm} p{2.6cm} p{2.6cm}@{}}
\toprule
Channels active & Commit architecture & Speculative architecture\\
\midrule
none (A1--A3 enforced) & \OursV{} $\equiv$ \Bplus{}: identical decisions, $\bar D=0$ (Prop.~\ref{prop:eq}) & same \\
C1 only & $\bar D(\OursV-\Bplus)\ge0$, second order (Prop.~\ref{prop:c1}) & same \\
C1+C2+C3 & \textbf{primary:} $\bar D(\OursV-\Bplus)>0$ (Prop.~\ref{prop:c2}) & \textbf{primary:} $\bar D(\OursV-\Bten)\le\lamC(1+w)$, i.e.\ joint gains at most the compute bound (Prop.~\ref{prop:spec}) \\
\bottomrule
\end{tabular}
\end{table}

